\documentclass[11pt]{amsart}
\usepackage{amssymb,hyperref}
\usepackage[margin=1.1in]{geometry}
\usepackage{fancyvrb}
\usepackage[no-math]{fontspec}
\setmonofont{DejaVu Sans Mono}[Scale=0.85]
\fvset{fontsize=\footnotesize}
% compile with lualatex (latexmk -lualatex)
\emergencystretch=3em\sloppy

\newtheorem{theorem}{Theorem}
\newtheorem{proposition}[theorem]{Proposition}
\theoremstyle{remark}
\newtheorem{remark}[theorem]{Remark}

\title{A remark on OpenAI family 170: the case $r(4,t)$}
\author{[AUTHOR NAME]}
\date{October 9, 2026}

\begin{document}
\begin{abstract}
OpenAI's manuscripts in family~170 of \texttt{github.com/openai/math}
state the sharp logarithmic exponent
$r(s,t)=t^{s-1}/(\log t)^{s-2+o(1)}$ only for $s\ge5$
($s=5$ in a companion paper, $s\ge6$ in the general paper, whose
construction theorem is stated for projective dimension $d\ge5$).
We observe that the accompanying Lean formalization already proves the
prime-field construction for every $d\ge3$, and that the library's
lower-bound transfer and upper bound need only $d\ge1$ and $s\ge3$.
Instantiating at $d=3$ gives, with a ten-line Lean file whose axioms are
only \texttt{propext}, \texttt{Classical.choice}, \texttt{Quot.sound},
\[
\frac{t^3}{(\log t)^{2+\varepsilon}}\le r(4,t)\le C\frac{t^3}{(\log t)^2}
\qquad(t\ge t_0(\varepsilon)),
\]
so $r(4,t)=t^3/(\log t)^{2+o(1)}$. The previous best lower bound we are
aware of is $\Omega(t^3/(\log t)^4)$ (Mattheus--Verstra\"ete). No new
mathematics is claimed: the construction and every proof are OpenAI's.
\end{abstract}
\maketitle

\section{Statement}
Let $r(s,t)$ be the least $N$ such that every simple graph on $N$
vertices has a clique of order $s$ or an independent set of order $t$.

\begin{theorem}\label{thm:main}
There is $C>0$ such that for every $\varepsilon>0$ and all sufficiently
large integers $t$,
\[
\frac{t^3}{(\log t)^{2+\varepsilon}}\le r(4,t)\le C\frac{t^3}{(\log t)^2},
\qquad\text{hence}\qquad
\lim_{t\to\infty}\frac{3\log t-\log r(4,t)}{\log\log t}=2 .
\]
\end{theorem}

This is exactly the statement \texttt{MainBounds 4 $\wedge$ MainLimit 4}
in the notation of the comparator file
\texttt{lean/ComparatorChallenges/SharpLogRamsey.lean}, whose public
theorem is stated for $6\le s$.

\paragraph{Before / after.}
\begin{center}
\small\begin{tabular}{lp{0.33\textwidth}p{0.4\textwidth}}
 & lower bound for $r(4,t)$ & source\\\hline
Before & $c\,t^3/(\log t)^4$ & Mattheus--Verstra\"ete~\cite{MV}; Brada\v{c}~\cite{Bradac} at $s=4$\\
After & $t^3/(\log t)^{2+\varepsilon}$, every $\varepsilon>0$ & OpenAI 170 Lean library at $d=3$ (this remark)\\
Upper & $C\,t^3/(\log t)^2$ & Ajtai--Koml\'os--Szemer\'edi~\cite{AKS}; see also~\cite{LRZ}; also in the 170 library
\end{tabular}
\end{center}

\section{Where the statement comes from}
In the general manuscript (\emph{Sharp Logarithmic Exponents for Fixed
Off-Diagonal Ramsey Numbers}, Theorem~1.2) the construction reads: for
$d\ge5$, $0<\eta<1/10$ and every sufficiently large prime $q$, with
$\sigma=\log q$, there is a $K_{d+1}$-free graph on
$\lfloor q^d\sigma\rfloor$ vertices with independence number
$<\lfloor q\sigma^{1+\eta}\rfloor$ (the random ordered incident-flag graph
of $\mathrm{PG}(d,q)$). In Lean this is the proposition
\texttt{LowerTransfer.PrimeConstruction d}. The library contains
\begin{Verbatim}
theorem actual_prime_construction (i : ℕ) : PrimeConstruction (i+3)
theorem prime_construction (d : ℕ) (hd : 5 ≤ d) : PrimeConstruction d
theorem all_epsilon_lower (d : ℕ) (hd : 1≤d) (hprime : PrimeConstruction d)
    (ε : ℝ) (hε : 0<ε) : ∀ᶠ t : ℕ in atTop,
      (t:ℝ)^d/(Real.log t)^(((d-1:ℕ):ℝ)+ε) ≤ (ramsey (d+1) t :ℝ)
theorem Upper.ramsey_upper {s : ℕ} (hs : 3 ≤ s) : ∃ C : ℝ, 0 < C ∧ ∃ t₀ : ℕ,
    ∀ t : ℕ, t₀ ≤ t → (ramsey s t : ℝ) ≤ C*(t:ℝ)^(s-1)/(Real.log (t:ℝ))^(s-2)
\end{Verbatim}
(\texttt{Construction/PrimeConstruction.lean},
\texttt{Bounds/Lower.lean}, \texttt{Bounds/SamplingParameter.lean}).
The restriction $d\ge5$ appears only in the wrapper
\texttt{prime\_construction}, which is itself proved from
\texttt{actual\_prime\_construction (d-3)}; the formal proof of
\texttt{actual\_prime\_construction} is uniform in $i\ge0$, i.e.\ in
$d\ge3$ (vector dimension $i+4$). The public theorem
\texttt{main\_of\_primeConstruction} assumes $6\le s$ but its proof uses
only $3\le s$ (upper bound) and $1\le s-1$ (transfer).

The manuscript itself is consistent with this: its geometric core is
carried out in projective dimensions $2$ and $3$, and higher dimensions
$j\ge4$ are reduced to these by repeated projection
(Section~2, Proposition ``Repeated projection''). Neither manuscript
gives a reason for excluding $s=4$; the companion $r(5,t)$ paper remarks
that one of its counting bounds ``also covers dimension three''.
We did not re-audit the manuscript at $d=3$; the evidence offered here is
the machine-checked proof.

\section{The Lean certificate}
\begin{Verbatim}
import OAI.Combinatorics.SharpRamsey.Main
open OAI.SharpLogRamsey

theorem r4_lower (ε : ℝ) (hε : 0 < ε) :
    ∀ᶠ t : ℕ in Filter.atTop,
      (t:ℝ)^3/(Real.log t)^(((3-1:ℕ):ℝ)+ε) ≤ (ramsey (3+1) t : ℝ) :=
  LowerTransfer.all_epsilon_lower 3 (by norm_num)
    (LowerTransfer.actual_prime_construction 0) ε hε

theorem main_four : MainBounds 4 ∧ MainLimit 4 := by
  have hb : MainBounds 4 := by
    obtain ⟨C, hC, t₀, ht₀⟩ := Upper.ramsey_upper (s := 4) (by norm_num)
    refine ⟨C, hC, fun ε hε => ?_⟩
    have hl := LowerTransfer.all_epsilon_lower 3 (by norm_num)
      (LowerTransfer.actual_prime_construction 0) ε hε
    obtain ⟨t₁, ht₁⟩ := Filter.eventually_atTop.mp hl
    exact ⟨max t₀ t₁, fun t ht => ⟨ht₁ t ((le_max_right _ _).trans ht),
      ht₀ t ((le_max_left _ _).trans ht)⟩⟩
  exact ⟨hb, mainLimit_of_mainBounds 4 hb⟩
\end{Verbatim}
Here \texttt{ramsey s t} is
$\inf\{N:\ \forall G\text{ on }\mathrm{Fin}\,N,\ G\text{ has an }s\text{-clique or }G^{c}\text{ has a }t\text{-clique}\}$,
which is $r(s,t)$; the definitions in \texttt{Model.lean} are textually
identical to those of the comparator file. Pinned to
\texttt{openai/math} commit \texttt{fd4aeeb2ee4f} (the
\texttt{SharpRamsey} directory is unchanged since the initial commit
\texttt{adc7f1241b42}), Lean \texttt{v4.34.1}, both theorems elaborate and
\texttt{\#print axioms} reports only \texttt{propext},
\texttt{Classical.choice}, \texttt{Quot.sound}. The \texttt{SharpRamsey}
sources contain no \texttt{sorry}, \texttt{axiom}, \texttt{native\_decide},
\texttt{implemented\_by}, \texttt{unsafe} or \texttt{set\_option}, and import
nothing from \texttt{OAI} outside that directory.

\begin{remark}
The same observation at $i=1$ ($d=4$) recovers the $r(5,t)$ statement
from the general library, independently of the separate
\texttt{RamseyFive} formalization. We do not pursue this.
\end{remark}

\section*{Scope}
Nothing here is claimed as new mathematics; the construction, its proof
and its formalization are OpenAI's. The remark only records that the
existing formal proof covers $s=4$, which the manuscripts and the
comparator statement do not state. To our knowledge (search of arXiv,
the web and the \texttt{openai/math} GitHub repository up to 2026-10-09)
no earlier source states a lower bound for $r(4,t)$ with logarithmic power
below~$4$.

\medskip\noindent\emph{Developed with AI assistance (Anthropic Claude,
Opus 5.5) under [AUTHOR NAME]'s direction; independent adversarial reviews
were run by separate AI sessions and are not human peer review.}

\begin{thebibliography}{9}
\bibitem{AKS} M.~Ajtai, J.~Koml\'os, E.~Szemer\'edi, A note on Ramsey
numbers, \emph{J. Combin. Theory Ser. A} \textbf{29} (1980), 354--360.
\bibitem{Bradac} D.~Brada\v{c}, Off-diagonal Ramsey numbers,
arXiv:2605.28793v3 (2026).
\bibitem{LRZ} Y.~Li, C.~C.~Rousseau, W.~Zang, Asymptotic upper bounds for
Ramsey functions, \emph{Graphs Combin.} \textbf{17} (2001), 123--128.
\bibitem{MV} S.~Mattheus, J.~Verstra\"ete, The asymptotics of $r(4,t)$,
\emph{Ann. of Math.} \textbf{199} (2024), 919--941.
\bibitem{OAI170} OpenAI, Sharp logarithmic exponents for fixed
off-diagonal Ramsey numbers; The sharp logarithmic exponent of $r(5,t)$,
OpenAI Math Release (2026), \url{https://github.com/openai/math}, family 170.
\end{thebibliography}
\end{document}
